% front/07_kor_abstract.tex

\clearpage
\addcontentsline{toc}{chapter}{\PlainTOCtext{국문초록}}

% ① "국문초록" 제목 (왼쪽, 16pt Bold)
{\fontsize{16}{20}\selectfont\bfseries 국문초록}\par
\vspace{15mm}  % 1.5cm

% ② 한글 논문 제목 (16pt Bold)
\begin{center}
  \begin{minipage}{0.9\textwidth}
    \centering
    {\fontsize{16}{20}\selectfont\bfseries
    관측성과 불확실성 기반 동적\par
    복구목표시간(DRTO) 프레임워크 연구
    }\par
  \end{minipage}
\end{center}
\vspace{15mm}  % 1.5cm

% ③ 이름/학과/대학원 – 오른쪽 정렬 + 2cm 안쪽
\begin{flushright}
  \hspace*{-2cm}
  {\fontsize{14}{22}\selectfont
    \setlength{\baselineskip}{22pt}
    이\ \ 창\ \ 호\par
    재난안전관리학과\par
    숭실대학교 대학원\par
  }
\end{flushright}

\vspace{15mm}  % 1.5cm

% ④ 본문 (11pt, 줄간 200%)
{\fontsize{11}{22}\selectfont
\setstretch{2.0}\justifying
\setlength{\parindent}{2em}

업무연속성관리(BCM)에서 복구목표시간(RTO)은 업무영향분석(BIA) 시점에 고정값으로 설정되어, 이후 운영 환경의 변화와 무관하게 유지된다. 이러한 정적 RTO 접근법은 시스템 관측성 수준의 동적 변동, 복구 과정에 내재된 불확실성의 확률적 특성, 사이버 공격이나 복합 재난 등 극단적 사건의 파레토 위험을 구조적으로 반영하지 못하는 한계를 지닌다.

본 연구는 이러한 한계를 극복하기 위해, \chg{관측성과 불확실성을 개별 시그모이드 바운딩 함수로 정식화한 수식 모델, 정적 기준선에서 파레토 불확실성까지의 4단계 운영 조건(C0--C3), 조건당 10,000회 몬테카를로 시뮬레이션, 효과 크기와 CDF 교차점 기반 분할 회귀 및 CVaR 꼬리 위험 정량화를 포함하는 다차원 비교 분석, 다중 시드($10^8$회) 강건성 검증, 기초 검증과 10개 연구가설 검증, 전문가 평가를 결합한 6단계 다층적 타당성 검증 체계(L0--L5), 6개 산업 시나리오 범용성 검증}으로 구성되는 DRTO 프레임워크를 제안한다. 핵심 수식은 $\DRTO = R_0 + E_{O,\mathrm{bnd}} + E_{U,\mathrm{bnd}}$로 정의되며, 관측성 효과 $E_{O,\mathrm{bnd}} = -R_0 \cdot S(z_O)$는 복구 시간을 기준선 $R_0$로부터 하향 조정하고, 불확실성 효과 $E_{U,\mathrm{bnd}} = (R_{\max} - R_0) \cdot S(z_U)$는 상향 조정한다. 수정 시그모이드 $S(z) = 2/(1 + e^{-z}) - 1$에 의해 $\DRTO$는 물리적 경계 $(0,\; R_{\max})$ 내에서 점근적으로 보장되며, $R_0 = 6.0$시간, $R_{\max} = 24.0$시간($= 4R_0$), $\kappa = 1.2$를 기본 설정으로 채택하였다. \fpara{$\kappa = 1.2$는 C1 조건의 격자 탐색에서 도출된 결과이며, 본 연구는 이를 \emph{환경별 최적값}이 아닌 BCM 일반 운영의 \emph{\chg{잠정적이고 보수적인} 실무 기본값}으로 채택한다. C2/C3 조건에서의 적합성은 가설 H2--H10 PASS의 통합 해석으로 \emph{운영 적합성을 시사}하며, 환경 무관 $\kappa$의 \emph{존재 여부 탐색}은 후속 연구의 과제로 남긴다.}

모델은 효율비 $\eta = \DRTO / R_0$를 기준으로 네 가지 운영 조건을 정의한다. \chg{정적 기준선인 C0는 $\eta = 1.000$이고, 관측성만 활성화한 C1은 $\eta \leq 1.000$이며, 가우시안 불확실성을 결합한 C2는 $\eta \geq \eta_{\mathrm{C1}}$, 파레토 불확실성의 C3는 꼬리 영역에서 $\eta \geq \eta_{\mathrm{C2}}$를 보인다.} 두 분포의 기대값을 $3.0$으로 통일하여 분포 형태의 차이만을 비교하도록 설계하였다.

조건당 $n = 10{,}000$ 몬테카를로 시뮬레이션 결과, 공통 경로에서 C0($1.000$) $\to$ C1($0.853$)으로 관측성에 의해 평균 $14.7\%$ 단축되었으며(Cohen's $d = -1.180$), C1에서 분기하여 C2($1.682$)와 C3($1.514$) 모두 불확실성 프리미엄에 의해 상향되었다. C2와 C3는 동일한 평균값($3.0$)을 갖도록 설계되었음에도 불구하고, 극단 백분위 영역에서는 현저한 차이를 보였다. P99 기준에서 C3의 $\DRTO$는 약 $13.2$시간으로, C2의 약 $8.4$시간보다 약 $57\%$ 높은 값을 나타냈으며, 이는 파레토 분포의 $U_{\mathrm{raw}}$ 99백분위가 가우시안 대비 약 $4.1$배($21.85$ vs. $5.33$) 큰 꼬리 특성에 기인한다.

두 불확실성 분포의 CDF가 교차하는 P85.8($\eta \approx 2.42$)을 기준으로, 핵심 영역(Core, $\leq$ P85.8)과 꼬리 영역(Tail, $>$ P85.8)으로 분할 회귀를 수행한 결과, 관측성 가중치 $k_O$의 회귀 계수가 Core($-0.798$)에서 Tail($-0.415$)로 $0.52$배 감쇠되는 이중채널 감쇠 효과를 발견하였다. 이는 극단적 불확실성 하에서 관측성의 복구 단축 효과가 구조적으로 약화됨을 의미한다. CVaR$_{99}$ 기준 C3가 C2 대비 $+24.0\%$ 높아, 파레토 위험이 극단 영역의 복구 부담을 불균형적으로 증폭시킴을 정량적으로 입증하였다.

본 프레임워크는 다층적 검증 체계를 통해 타당성을 확보하였다. 10개 \m{연구 가설}(H1--H10)이 모두 채택되었으며, BCM/DR 분야 전문가 5명의 정량적 설문은 전체 평균 $4.46/5.0$점($p = 0.011$, Cohen's $d = 1.65$)과 높은 내적 일관성(Cronbach's $\alpha = 0.89$, S-CVI/Ave$\,= 0.985$)을 보였고, 심층 인터뷰에서는 5대 주제 15개 하위 주제 모두에서 $80\%$ 이상 동의($100\%$ 충족)를 확인하였다. 2차 회귀 분석은 C1에서 $R^2 = 1.000$, C2에서 $R^2 = 0.998$의 사실상 완전한 설명력을 달성하였고, C3는 전체 $R^2 = 0.858$이나 P85.8 분할 시 Core $R^2 = 0.999$, Tail $R^2 = 0.710$으로서 분할 회귀의 유효성을 입증하였다. 또한 6개 산업 시나리오(\chg{$R_0 = 0.5$--$8.0$시간})에서 P99 비율이 평균 $1.57$배로 일관되어 모델의 범용 적용 가능성을 확인하였다.

본 연구는 관측성의 이점과 불확실성의 패널티를 개별 바운딩으로 통합하고, 특히 복구 계획에 대한 Heavy-Tail 분포의 불균형적 영향과 이중채널 감쇠 메커니즘을 규명함으로써, \chg{통계적 근거에 기반한 DRTO 조정 방법론과 더불어, 회귀식 기반 $\DRTO$ 산출 도구와, 이중채널 감쇠에 기반한 자원 배분 원칙 및 스트레스 테스트 설계 기준을} 업무연속성관리 분야에 기여한다. \fpara{본 연구의 학술적 한계로서, $\kappa = 1.2$의 환경 무관 보편 최적성은 본 연구에서 입증되지 않으며 이는 후속 연구의 탐색 과제로 남긴다.}

\vspace{10mm}
\noindent\textbf{주제어:} 복구목표시간(RTO), 동적복구목표시간(DRTO), 업무연속성관리(BCM), 관측성, 불확실성정량화, 시그모이드바운딩, 중꼬리분포

}
